% !PRINTED 671
\weberpage{671}
Es sei also $\Omega$ ein Körper $n^{\text{ten}}$ Grades, und

\weq{1}{\Omega_1,\ \Omega_2,\ \ldots,\ \Omega_n}

die conjugirten Körper. Die irreducible Gleichung $n^{\text{ten}}$ Grades, deren Wurzeln uns diese conjugirten Körper bestimmen, wird im Allgemeinen sowohl reelle als imaginäre Wurzeln haben, von denen aber die letzteren immer paarweise vorkommen, so dass zu jedem imaginären Körper $\Omega^{(i)}$ ein anderer gehört, dessen Zahlen mit denen von $\Omega^{(i)}$ conjugirt imaginär sind, der also aus $\Omega^{(i)}$ durch die Substitution $(i, -i)$ hervorgeht.

Solche imaginäre Paare fassen wir zu einer Einheit zusammen und bezeichnen \gesperrt{die Anzahl der reellen Körper und der Paare imaginärer Körper der Reihe (1)} mit $\nu$. Sind alle conjugirten Körper reell, so ist $n = \nu$; sind sie alle imaginär, so ist $n = 2\nu$.

Ausser im Falle des rationalen Körpers ist $\nu$ nur noch in dem Falle eines imaginären quadratischen Körpers $= 1$. Sonst ist $\nu$ immer grösser als 1.

Die Anzahl der imaginären Paare ist $n - \nu$, und die Anzahl der reellen Körper $2\nu - n$.

In der Determinante, durch die in \S.\,145 die Quadratwurzel aus der Grundzahl, $\sqrt{\Delta}$, definirt ist, vertauschen sich durch die Substitution $(i, -i)$ je zwei conjugirt imaginäre Reihen, und $\sqrt{\Delta}$ wechselt also sein Zeichen dann, wenn diese Zahl ungerade ist, sonst nicht. Im ersten Falle ist $\sqrt{\Delta}$ imaginär, im zweiten reell. Daraus ergiebt sich also, dass die \gesperrt{Grundzahl $\Delta$ des Körpers $\Omega$ positiv oder negativ ist, je nachdem $n - \nu$ gerade oder ungerade ist.}

Behalten wir von zwei conjugirt imaginären Körpern nur den einen bei, so ergiebt sich die Reihe der conjugirten Körper

\[ \Omega_1,\ \Omega_2,\ \ldots,\ \Omega_\nu, \]

denen wir ein Zeichensystem

\[ \delta_1,\ \delta_2,\ \ldots,\ \delta_\nu \]

mit der Bestimmung zuordnen, dass $\delta_s = 1$ sein soll, wenn $\Omega_s$ reell, und $= 2$, wenn $\Omega_s$ imaginär ist, also $\Sigma \delta_s = n$.

Es möge $\eta$ eine von Null verschiedene Zahl des Körpers $\Omega$ bedeuten und

\[ \eta_1,\ \eta_2,\ \ldots,\ \eta_n \]

die conjugirten Zahlen. Diesem Zahlensysteme ordnen wir ein anderes Zahlensystem zu

\[ \lambda_1,\ \lambda_2,\ \ldots,\ \lambda_n, \]
%CONT>>
